% ==============================================================================
% CHAPTER 2, PART B: SECTIONS 2.4 - 2.7
% Emergent Gravity, Quantum Nature of the Equivalence Principle, and Horizon Thermodynamics
% Translation from Ukrainian to English (PRD / JHEP style)
% Source line range: 2640 - 3246
% Equations: (2.26) - (2.71) [linear sequence (2.76) through (2.134)]
% ==============================================================================

\subsection{Vacuum Phase Density Gradient, Margolus-Levitin Theorem, and Quantum Refraction of de Broglie Waves}\label{sec:2.4}

\subsubsection{Vacuum Phase Density Gradient and the Margolus-Levitin Theorem \cite{ref16}}\label{sec:2.4.1}

In classical Newtonian mechanics and Einstein's General Relativity, the equivalence principle between inertial and gravitational mass ($m_{\text{in}} \equiv m_{\text{grav}}$) is postulated as a primary empirical fact without derivation from deeper quantum foundations. Within the framework of contour phase geometrodynamics developed in this work, this principle receives a rigorous quantum-kinematic proof without circular reliance on Einstein's metric field equations.

Consider the local progression of quantum processes in a system situated within a static gravitational field possessing a timelike Killing vector $\xi = \partial_t$. In the contour phase approach, a stationary quantum state with intrinsic Compton phase frequency $\omega_0$ acquires a phase accumulation along its worldline generated by the canonical 1-form $\mathbf{\Pi}_\mu dx^\mu = -\hbar \omega_0 d\tau$, where $\tau$ is the proper time along the trajectory.

In a static field, the relation between the proper time differential $d\tau$ and the time of an infinitely distant observer (coordinate time $t$) is given by the lapse function $N(\mathbf{r}) = \sqrt{-g_{00}(\mathbf{r})}$:

\begin{equation}\label{eq:2.26}\label{eq:2.76}\tag{2.26}
d\tau = \sqrt{-g_{00}(\mathbf{r})} \, dt = N(\mathbf{r}) \, dt = \sqrt{1 + \frac{2\Phi(\mathbf{r})}{c^2}} \, dt \approx \left( 1 + \frac{\Phi(\mathbf{r})}{c^2} \right) dt = \left( 1 - \frac{GM}{r c^2} \right) dt,
\end{equation}

where $G = 6.67430 \cdot 10^{-11}\text{ m}^3\cdot\text{kg}^{-1}\cdot\text{s}^{-2}$ is Newton's gravitational constant;  
$M$ is the mass of the gravitating source;  
$r$ is the radial distance;  
$c = 2.99792458 \cdot 10^8\text{ m}\cdot\text{s}^{-1}$ is the speed of light in vacuum.

The generator of time evolution along the coordinate time $t$, namely the coordinate quantum energy of the state $\mathcal{E}_{\text{coord}}(\mathbf{r})$, is determined by the projection of the canonical phase momentum onto the timelike Killing vector $\partial_t$:

\begin{equation}\label{eq:2.27}\label{eq:2.77}\tag{2.27}
\mathcal{E}_{\text{coord}}(\mathbf{r}) = i\hbar \frac{\partial}{\partial t} = N(\mathbf{r}) \, \mathcal{E}_0 = \hbar \omega_0 \sqrt{1 + \frac{2\Phi(\mathbf{r})}{c^2}} \approx \hbar \omega_0 \left( 1 - \frac{GM}{r c^2} \right),
\end{equation}

where $\mathcal{E}_0 = \hbar \omega_0$ is the unperturbed proper energy of the quantum state at infinity. It must be emphasized that in Eq.~\eqref{eq:2.27}, the energy shift arises purely geometrically through the lapse function $N(\mathbf{r})$ without an a priori introduction of a ``gravitational mass.''

According to the fundamental \textbf{Margolus-Levitin theorem} (1998) \cite{ref16}, the minimum coordinate time interval $\Delta t_{\min}$ required for a quantum system to evolve between two mutually orthogonal states (the period of a quantum clock) is strictly bounded by the average coordinate energy of the system $\mathcal{E}_{\text{coord}}(\mathbf{r})$:

\begin{equation}\label{eq:2.28}\label{eq:2.78}\tag{2.28}
\Delta t_{\min}(\mathbf{r}) = \frac{\pi \hbar}{2 \mathcal{E}_{\text{coord}}(\mathbf{r})} = \frac{\pi \hbar}{2 \hbar \omega_0 \left( 1 - \frac{GM}{r c^2} \right)} \approx \Delta t_{\min, 0} \left( 1 + \frac{GM}{r c^2} \right),
\end{equation}

where $\Delta t_{\min}$ is the fundamental step of discrete quantum evolution;  
$\hbar = 1.054571817 \cdot 10^{-34}\text{ J}\cdot\text{s}$ is the reduced Planck constant.

Accordingly, the local rate frequency of quantum phase transitions of the clock, $\nu(r) = 1/\Delta t_{\min}(r)$, decreases monotonically with proximity to the massive center:

\begin{equation}\label{eq:2.29}\label{eq:2.79}\tag{2.29}
\nu(r) = \nu_0 \left( 1 - \frac{GM}{r c^2} \right),
\end{equation}

where $\nu_0 = 2 \omega_0 / \pi$ is the unperturbed vacuum transition frequency. Relation~\eqref{eq:2.29} derives the phenomenon of gravitational time dilation directly from the quantum kinematics of the Margolus-Levitin theorem and the temporal projection of the contour phase, completely eliminating circular dependence on the initial postulation of $m_{\text{grav}}$.

The reduction in the rate of quantum transitions~\eqref{eq:2.29} generates a radial spatial gradient of the vacuum phase density: the vacuum in a gravitational field behaves as a medium with a spatially inhomogeneous phase accumulation index.

\subsubsection{Quantum Refraction of de Broglie Wave Packets}\label{sec:2.4.2}

Consider the wave dynamics of a de Broglie quantum wave packet in an inhomogeneous vacuum phase medium. The wave function of the relativistic field satisfies the generalized wave equation derived from the master-operator:

\begin{equation}\label{eq:2.30}\label{eq:2.80}\tag{2.30}
\left( g^{\mu\nu}\nabla_\mu \nabla_\nu + \frac{m^2 c^2}{\hbar^2} \right) \psi(x) = 0.
\end{equation}

In the eikonal approximation of geometric optics, $\psi(x) = \mathcal{A}(x) \exp(i S(x)/\hbar)$, where the 4-wavevector is $k_\mu = -\partial_\mu S / \hbar = (\omega/c, -\mathbf{k})$, Eq.~\eqref{eq:2.30} transforms into the fundamental dispersion relation:

\begin{equation}\label{eq:2.31}\label{eq:2.81}\tag{2.31}
g^{\mu\nu} k_\mu k_\nu = \frac{m^2 c^2}{\hbar^2} \implies g^{00} \frac{\omega^2}{c^2} + g^{ij} k_i k_j + \frac{m^2 c^2}{\hbar^2} = 0.
\end{equation}

In the weak field of a static center, the isotropic components of the contravariant metric tensor are $g^{00} \approx -(1 - 2\Phi/c^2)$ and $g^{ij} \approx (1 + 2\Phi/c^2)\delta^{ij}$. Substituting them into the dispersion relation~\eqref{eq:2.31}:

\begin{equation}\label{eq:2.32}\label{eq:2.82}\tag{2.32}
\frac{\omega^2}{c^2} \left( 1 - \frac{2\Phi(\mathbf{r})}{c^2} \right) - k^2 \left( 1 + \frac{2\Phi(\mathbf{r})}{c^2} \right) - \frac{m^2 c^2}{\hbar^2} = 0.
\end{equation}

Expressing the local angular frequency of the wave packet $\omega(\mathbf{k}, \mathbf{r})$:

\begin{equation}\label{eq:2.33}\label{eq:2.83}\tag{2.33}
\omega(\mathbf{k}, \mathbf{r}) = c \sqrt{ \left( 1 + \frac{2\Phi}{c^2} \right) \left[ \frac{m^2 c^2}{\hbar^2} + k^2 \left( 1 + \frac{2\Phi}{c^2} \right) \right] } \approx c \sqrt{ \frac{m^2 c^2}{\hbar^2}\left(1 + \frac{2\Phi}{c^2}\right) + k^2 \left(1 + \frac{4\Phi}{c^2}\right) }.
\end{equation}

The group velocity of energy transport by the de Broglie wave packet is:

\begin{equation}\label{eq:2.34}\label{eq:2.84}\tag{2.34}
\mathbf{v}_g = \nabla_{\mathbf{k}} \omega = \frac{c^2 \mathbf{k}}{\omega} \left( 1 + \frac{4\Phi}{c^2} \right) \approx \frac{\hbar \mathbf{k}}{m} = \mathbf{v}.
\end{equation}

From the wave kinematics of the Hamilton-Jacobi eikonal, the spatial gradient of frequency $\nabla_\mathbf{r} \omega$ governs the time evolution of the wavevector:

\begin{equation}\label{eq:2.35}\label{eq:2.85}\tag{2.35}
\frac{d\mathbf{k}}{dt} = -\nabla_{\mathbf{r}} \omega(\mathbf{k}, \mathbf{r}).
\end{equation}

Differentiating the dispersion relation~\eqref{eq:2.33} with respect to the spatial coordinates $\mathbf{r}$, we obtain:

\begin{equation}\label{eq:2.36}\label{eq:2.86}\tag{2.36}
\nabla_{\mathbf{r}}\omega = \frac{c^2}{2\omega} \left[ \frac{m^2 c^2}{\hbar^2} \frac{2\nabla\Phi}{c^2} + k^2 \frac{4\nabla\Phi}{c^2} \right] = \frac{m^2 c^2}{\hbar^2 \omega} \nabla\Phi + \frac{2 k^2}{\omega} \nabla\Phi.
\end{equation}

From this, the deep physical nature of the distinction between slow massive matter and light is revealed:

\begin{enumerate}
\item \textbf{For non-relativistic massive quantum particles ($v \ll c$):}  
In this regime, $\hbar \omega \approx m c^2 \gg \hbar c k$, and the ratio of the wavenumber term to the mass term is $(v/c)^2 \ll 1$. The second term in Eq.~\eqref{eq:2.36}, originating from the spatial metric deformation, is suppressed by the factor $(v/c)^2 \to 0$. Accordingly:

\begin{equation}\label{eq:2.37}\label{eq:2.87}\tag{2.37}
\nabla_{\mathbf{r}}\omega \approx \frac{m}{\hbar} \nabla\Phi(\mathbf{r}).
\end{equation}

Substituting Eq.~\eqref{eq:2.37} into the eikonal wave equation~\eqref{eq:2.35} and multiplying by $\hbar$:

\begin{equation}\label{eq:2.38}\label{eq:2.88}\tag{2.38}
\frac{d(\hbar \mathbf{k})}{dt} = m \frac{d\mathbf{v}}{dt} = m \mathbf{a} = -m \nabla \Phi(\mathbf{r}) = -m \frac{GM}{r^2}\hat{\mathbf{r}}.
\end{equation}

Cancelling the scalar mass $m$, we obtain the universal free-fall acceleration:

\begin{equation}\label{eq:2.39}\label{eq:2.89}\tag{2.39}
\mathbf{a} = -\nabla\Phi(\mathbf{r}) = -\frac{GM}{r^2} \hat{\mathbf{r}}.
\end{equation}

Thus, for massive bodies, the spatial curvature is suppressed kinematically ($\mathcal{O}(v^2/c^2)$), and the acceleration is determined exclusively by the gradient of the temporal lapse function.

\item \textbf{For massless electromagnetic waves / photons ($m = 0$, $v = c$):}  
For photons $m = 0$, so $\omega = c k \left(1 + \frac{2\Phi}{c^2}\right)^{1/2}\left(1 - \frac{2\Phi}{c^2}\right)^{-1/2} \approx c k \left(1 + \frac{2\Phi}{c^2}\right)$. Both contributions---temporal ($1 \times \Phi/c^2$) and spatial ($1 \times \Phi/c^2$)---enter with equal unit weight, since $v = c$:

\begin{equation}\label{eq:2.40}\label{eq:2.90}\tag{2.40}
v_{\text{light}}(r) = \frac{\omega}{k} \approx c \left( 1 + \frac{2\Phi(r)}{c^2} \right) = c \left( 1 - \frac{2GM}{r c^2} \right).
\end{equation}

Accordingly, the effective optical refractive index of the vacuum for electromagnetic waves is:

\begin{equation}\label{eq:2.41}\label{eq:2.91}\tag{2.41}
n_{\text{opt}}(r) = \frac{c}{v_{\text{light}}(r)} \approx 1 - \frac{2\Phi(r)}{c^2} = 1 + \frac{2GM}{r c^2}.
\end{equation}
\end{enumerate}

The factor of 2 is a direct mathematical consequence of the equal participation of spatial and temporal components of the phase metric perturbation at $v = c$. This completely explains the doubled deflection angle of light by the Sun, $\theta = 4GM/(bc^2)$, while simultaneously preserving the standard Newtonian acceleration~\eqref{eq:2.39} for slow quantum bodies.

\subsection{Analytical Proof of the Equivalence Principle ($m_{\text{in}} \equiv m_{\text{grav}}$) and COW Experiment}\label{sec:2.5}

\subsubsection{Rigorous Analytical Proof of the Identity of Masses}\label{sec:2.5.1}

We formulate and prove the fundamental theorem of this research.

\begin{theorem}[Quantum-Phase Origin of the Equivalence Principle]\label{thm:wep}
The universal wave functional of the Universe $\mathbf{\Psi}[\Sigma]$ on the space of closed contours satisfies the Wheeler-DeWitt superconstraint $\hat{\mathcal{H}}_\perp \mathbf{\Psi} \equiv 0$, in which the particle mass $\mathcal{M}_0$ emerges as the unique invariant scalar eigenvalue of the Casimir operator of the bundle isometry group:
\begin{equation*}
\hat{P}_\mu \hat{P}^\mu = \mathcal{M}_0^2 c^2.
\end{equation*}
Due to the unity of the phase action in the four-dimensional superspace of 3-metrics and matter fields, the inertial mass determining the resistance to acceleration and the gravitational mass determining the interaction with the vacuum phase gradient are identically equal:
\begin{equation}\label{eq:2.42}\label{eq:2.92}\tag{2.42}
m_{\text{in}} \equiv m_{\text{grav}} \equiv \mathcal{M}_0,
\end{equation}
and the free-fall acceleration of any particle is strictly universal and independent of its mass, internal quantum structure, and energy spectrum:
\begin{equation}\label{eq:2.43}\label{eq:2.93}\tag{2.43}
\mathbf{a} = -\frac{GM}{r^2} \hat{\mathbf{r}} = \mathbf{g}(\mathbf{r}).
\end{equation}
\end{theorem}

\begin{proof}
Consider the non-perturbative action functional of a material Feynman phase loop in the spacetime representation $(\mathcal{M}, g_{\mu\nu})$. For a quasiclassical wave packet forming a localized particle, the action along the worldline $\gamma$ is given by the integral:

\begin{equation}\label{eq:2.44}\label{eq:2.94}\tag{2.44}
\mathcal{S}[\gamma] = -\mathcal{M}_0 c \int_{\gamma} d\tau = -\mathcal{M}_0 c \int \sqrt{-g_{\mu\nu} \dot{x}^\mu \dot{x}^\nu} \, d\lambda,
\end{equation}

where $\mathcal{M}_0$ is the invariant Casimir scalar mass, and $\lambda$ is a monotonic parameter along the contour.

We apply the canonical $3+1$ Arnowitt-Deser-Misner (ADM) decomposition to the spacetime metric with lapse function $N = \sqrt{-g_{00}}$, vanishing shift ($N^i = 0$), and spatial 3-metric $h_{ij}$:

\begin{equation}\label{eq:2.45}\label{eq:2.95}\tag{2.45}
-g_{\mu\nu}\dot{x}^\mu \dot{x}^\nu = N^2 c^2 \dot{t}^2 - h_{ij}\dot{x}^i \dot{x}^j.
\end{equation}

Choosing coordinate time as the trajectory parameter ($\lambda = t, \dot{t} = 1$), the Lagrangian of the particle takes the form:

\begin{equation}\label{eq:2.46}\label{eq:2.96}\tag{2.46}
L(x, \mathbf{v}) = -\mathcal{M}_0 c \sqrt{N^2(\mathbf{r}) c^2 - h_{ij}\dot{x}^i \dot{x}^j} = -\mathcal{M}_0 c^2 N(\mathbf{r}) \sqrt{1 - \frac{h_{ij}\dot{x}^i \dot{x}^j}{c^2 N^2(\mathbf{r})}}.
\end{equation}

We expand the Lagrangian~\eqref{eq:2.46} in powers of the small non-relativistic parameter $(v/c)^2 \ll 1$ in the weak-field regime $N(\mathbf{r}) \approx 1 + \Phi(\mathbf{r})/c^2$:

\begin{equation}\label{eq:2.47}\label{eq:2.97}\tag{2.47}
L(x, \mathbf{v}) \approx -\mathcal{M}_0 c^2 \left( 1 + \frac{\Phi(\mathbf{r})}{c^2} \right) \left[ 1 - \frac{1}{2}\frac{v^2}{c^2} \right] \approx -\mathcal{M}_0 c^2 + \frac{1}{2}\mathcal{M}_0 v^2 - \mathcal{M}_0 \Phi(\mathbf{r}).
\end{equation}

From the resulting Lagrangian~\eqref{eq:2.47}, the definitions of canonical momentum and generalized force follow directly:

\begin{enumerate}
\item \textbf{Canonical momentum (measure of system inertia):}

\begin{equation}\label{eq:2.48}\label{eq:2.98}\tag{2.48}
p_i = \frac{\partial L}{\partial v^i} = \mathcal{M}_0 v_i \implies m_{\text{in}} \equiv \frac{\partial p_i}{\partial v_i} = \mathcal{M}_0.
\end{equation}

\item \textbf{Gravitational force (measure of interaction with the vacuum phase gradient):}

\begin{equation}\label{eq:2.49}\label{eq:2.99}\tag{2.49}
F_i^{\text{grav}} = \frac{\partial L}{\partial x^i} = -\mathcal{M}_0 \nabla_i \Phi(\mathbf{r}) \implies m_{\text{grav}} \equiv \frac{\partial F_i^{\text{grav}}}{\partial (-\nabla_i \Phi)} = \mathcal{M}_0.
\end{equation}
\end{enumerate}

Because both the kinetic term $\frac{1}{2}\mathcal{M}_0 v^2$ and the potential term $-\mathcal{M}_0 \Phi(\mathbf{r})$ originate from the series expansion of the very same square root of invariant worldline length with a single scalar invariant $\mathcal{M}_0$, splitting mass into two independent entities $m_{\text{in}}$ and $m_{\text{grav}}$ is mathematically impossible within the non-perturbative DeWitt supermetric:

\begin{equation}\label{eq:2.50}\label{eq:2.100}\tag{2.50}
\frac{m_{\text{grav}}}{m_{\text{in}}} \equiv \frac{\mathcal{M}_0}{\mathcal{M}_0} \equiv 1 \implies m_{\text{in}} \equiv m_{\text{grav}}.
\end{equation}

Substituting Eqs.~\eqref{eq:2.48} and \eqref{eq:2.49} into the Euler-Lagrange equations of motion $\frac{d p_i}{dt} - F_i^{\text{grav}} = 0$:

\begin{equation}\label{eq:2.51}\label{eq:2.101}\tag{2.51}
\mathcal{M}_0 \frac{d\mathbf{v}}{dt} = -\mathcal{M}_0 \nabla \Phi(\mathbf{r}) \implies \mathbf{a} = \frac{d\mathbf{v}}{dt} = -\nabla\Phi(\mathbf{r}) = -\frac{GM}{r^2}\hat{\mathbf{r}}.
\end{equation}

The scalar mass $\mathcal{M}_0$ strictly cancels identically on both sides of the dynamical equation. Consequently, Einstein's equivalence principle is a direct mathematical consequence of fiber bundle geometry and the uniqueness of the Casimir invariant of the contour functional $\mathbf{\Psi}$, completely depriving it of heuristic or postulational status. This completes the proof.
\end{proof}

\subsubsection{The Colella-Overhauser-Werner (COW) Quantum Experiment}\label{sec:2.5.2}

The most direct experimental confirmation of the wave-phase nature of the gravitational interaction is the celebrated Colella-Overhauser-Werner quantum experiment (COW experiment \cite{ref17}, 1975).

In the COW experiment, a monochromatic beam of thermal neutrons is split by a single-crystal silicon diffraction interferometer into two coherent, spatially separated beams that form a closed rectangular contour of height $H$ and length $L$ in the gravitational field of the Earth.

The potential difference between the upper and lower arms of the interferometer is:

\begin{equation}\label{eq:2.52}\label{eq:2.102}\tag{2.52}
\Delta \Phi = g H \sin\alpha,
\end{equation}

where $g \approx 9.81\text{ m}\cdot\text{s}^{-2}$ is the gravitational acceleration at the surface of the Earth;  
$H$ is the height of the rectangular contour;  
$\alpha$ is the tilt angle of the interferometer plane relative to the horizontal Earth surface.

The phase difference of the Feynman contour integral along the upper path $\gamma_1$ and the lower path $\gamma_2$ is:

\begin{equation}\label{eq:2.53}\label{eq:2.103}\tag{2.53}
\Delta \Phi_{\text{grav}} = \frac{1}{\hbar} \oint_\gamma p_\mu dx^\mu = \frac{1}{\hbar} \left( \int_{\gamma_1} p_1 dx - \int_{\gamma_2} p_2 dx \right).
\end{equation}

Employing the momentum expansion $p(z) = \sqrt{2m(E - mgz)} \approx p_0 - \frac{m^2 g z}{p_0}$:

\begin{equation}\label{eq:2.54}\label{eq:2.104}\tag{2.54}
\Delta \Phi_{\text{grav}} = -\frac{m^2 g H L \sin\alpha}{\hbar p_0} = -\frac{m^2 g \mathcal{A} \lambda_n \sin\alpha}{2\pi \hbar^2},
\end{equation}

where $\mathcal{A} = H \cdot L$ is the area of the interferometer contour;  
$\lambda_n = \frac{2\pi \hbar}{p_0}$ is the de Broglie wavelength of the thermal neutrons ($\lambda_n \sim 0.14\text{ nm}$);  
$m$ is the neutron mass ($m_n \approx 1.6749 \cdot 10^{-27}\text{ kg}$).

Experimental measurements of the interference fringe shift upon rotating the crystal through an angle $\alpha \in [-90^\circ, +90^\circ]$ confirmed the theoretical formula~\eqref{eq:2.54} with an accuracy better than $0.1\%$.

This provides direct experimental proof that \textbf{gravity in the quantum microworld operates directly as a quantum phase shift of the Feynman path integral}.

\subsubsection{Rigorous Fundamental Demarcation between the Weak (WEP) and Strong (SEP) Equivalence Principles}\label{sec:2.5.4}

From the perspective of theoretical physics and observational astrophysics, the formulated quantum contour framework establishes a clear demarcation line between different levels of the equivalence principle:

\begin{enumerate}
\item \textbf{Weak Equivalence Principle (WEP):}
Holds \textbf{strictly and universally} across all energy and spatial scales. Because inertial and gravitational masses are identical ($m_{\text{in}} \equiv m_{\text{grav}} \equiv \mathcal{M}_0$), the mass of a test body cancels identically in the canonical Euler-Lagrange equations:

\begin{equation}\label{eq:2.55a}\label{eq:2.105}\tag{2.55a}
\frac{d}{dt}\frac{\partial L}{\partial \dot{\mathbf{x}}} - \frac{\partial L}{\partial \mathbf{x}} = 0 \implies \ddot{x}^\mu + \Gamma^\mu_{\alpha\beta} \dot{x}^\alpha \dot{x}^\beta = 0.
\end{equation}

The trajectory of any quantum test body in the gravitational phase field is a pure geometric geodesic of the manifold, entirely independent of its mass, spin, baryon number, or internal quark structure. This has been confirmed by E\"otv\"os-type experiments, the MICROSCOPE satellite mission with a precision of $10^{-15}$, and COW neutron interferometry.

\item \textbf{Strong Equivalence Principle (SEP):}
Holds in local strong and Newtonian regimes ($g \gg a_0$), but is \textbf{violated at infrared acceleration scales ($g \lesssim a_0 \approx 1.2 \cdot 10^{-10}\text{ m/s}^2$)}. 
This violation is driven by the non-linearity of boundary conditions for the closed Feynman contour in an external gravitational field $\mathbf{g}_{\text{ext}}$ (\textbf{External Field Effect, EFE}). 

As a consequence of the EFE, the phase accumulation on galactic scales is governed by the total phase gradient:

\begin{equation}\label{eq:2.55b}\label{eq:2.106}\tag{2.55b}
\mathbf{g}_{\text{tot}} = \mathbf{g}_{\text{int}} + \mathbf{g}_{\text{ext}}.
\end{equation}

The internal dynamics of a gravitationally bound subsystem (for instance, a dwarf galaxy embedded in the field of a cluster or the outskirts of disk galaxies) ceases to be autonomous and depends explicitly on the strength of the uniform external field $\mathbf{g}_{\text{ext}}$, which in standard General Relativity is forbidden by Einstein's local elevator principle.
\end{enumerate}

\subsection{Microscopic Information-Combinatorial Derivation of Bekenstein-Hawking Black Hole Entropy}\label{sec:2.6}

\subsubsection{Event Horizon Thermodynamics and Derivation of the Bekenstein-Hawking Formula}\label{sec:2.6.1}

In 1972--1973, Jacob Bekenstein formulated the fundamental hypothesis that a black hole possesses a macroscopic thermodynamic entropy proportional to the area of its event horizon $A$ \cite{ref19}. In 1974--1975, Stephen Hawking established via quantum field theory in curved spacetime the celebrated numerical coefficient of $1/4$ \cite{ref18}:

\begin{equation}\label{eq:2.56}\label{eq:2.107}\tag{2.56}
S_{\text{BH}} = \frac{k_B c^3 A}{4 G \hbar} = \frac{k_B A}{4 l_P^2},
\end{equation}

where $S_{\text{BH}}$ is the thermodynamic entropy of the black hole;  
$k_B = 1.380649 \cdot 10^{-23}\text{ J}\cdot\text{K}^{-1}$ is the Boltzmann constant;  
$c = 2.99792458 \cdot 10^8\text{ m}\cdot\text{s}^{-1}$ is the speed of light in vacuum;  
$A$ is the area of the 2-dimensional event horizon surface of the black hole;  
$G = 6.67430 \cdot 10^{-11}\text{ m}^3\cdot\text{kg}^{-1}\cdot\text{s}^{-2}$ is Newton's gravitational constant;  
$\hbar = 1.054571817 \cdot 10^{-34}\text{ J}\cdot\text{s}$ is the reduced Planck constant;  
$l_P = \sqrt{\hbar G / c^3} \approx 1.616255 \cdot 10^{-35}\text{ m}$ is the fundamental Planck length.

Unlike loop quantum gravity, where the factor $1/4$ is achieved through ad hoc tuning of the free Barbero-Immirzi parameter $\gamma_{\text{BI}}$, in phase contour theory this coefficient is a rigorous analytical consequence of the First Law of horizon thermodynamics and the Unruh effect without any adjustable parameters.

Consider a static Schwarzschild black hole of mass $M$. The event horizon radius and the area of the 2-sphere are given by:

\begin{equation}\label{eq:2.57}\label{eq:2.108}\tag{2.57}
r_s = \frac{2GM}{c^2}, \quad A = 4\pi r_s^2 = \frac{16\pi G^2}{c^4} M^2.
\end{equation}

The surface gravity at the horizon is:

\begin{equation}\label{eq:2.58}\label{eq:2.109}\tag{2.58}
\kappa = \frac{c^4}{4GM} = \frac{c^2}{2 r_s},
\end{equation}

where $\kappa$ is the event horizon surface gravity.

According to the Unruh quantum effect, a uniformly accelerating observer near the horizon detects thermal vacuum radiation at the Hawking temperature \cite{ref18}:

\begin{equation}\label{eq:2.59}\label{eq:2.110}\tag{2.59}
T_H = \frac{\hbar \kappa}{2\pi k_B c} = \frac{\hbar c^3}{8\pi G M k_B}.
\end{equation}

We apply the First Law of horizon thermodynamics (the Clausius thermodynamic relation $dS = \delta Q / T_H$) for the quasistatic absorption of an elementary mass-energy $dE = c^2 dM$:

\begin{equation}\label{eq:2.60}\label{eq:2.111}\tag{2.60}
dS = \frac{c^2 dM}{T_H} = \frac{c^2 dM}{\frac{\hbar c^3}{8\pi G M k_B}} = \frac{8\pi G k_B}{\hbar c} M dM.
\end{equation}

Differentiating the expression for the horizon area $A(M)$~\eqref{eq:2.57}:

\begin{equation}\label{eq:2.61}\label{eq:2.112}\tag{2.61}
dA = \frac{32\pi G^2}{c^4} M dM \implies M dM = \frac{c^4}{32\pi G^2} dA.
\end{equation}

Substituting the differential $M dM$~\eqref{eq:2.61} directly into the thermodynamic relation~\eqref{eq:2.60}:

\begin{equation}\label{eq:2.62}\label{eq:2.113}\tag{2.62}
dS = \frac{8\pi G k_B}{\hbar c} \cdot \left( \frac{c^4}{32\pi G^2} dA \right) = \frac{k_B c^3}{4 \hbar G} dA = \frac{k_B}{4 l_P^2} dA.
\end{equation}

Integrating Eq.~\eqref{eq:2.62} from zero area to the macroscopic value $A$:

\begin{equation}\label{eq:2.63}\label{eq:2.114}\tag{2.63}
S_{\text{BH}} = \int_0^A \frac{k_B}{4 l_P^2} dA' = \frac{k_B A}{4 l_P^2}.
\end{equation}

The Bekenstein-Hawking formula is thus derived strictly analytically from first quantum-thermodynamic principles without invoking any arbitrary or adjustable constants.

\subsubsection{Quanta of Area per Bit of Information}\label{sec:2.6.2}

We now elucidate the microscopic statistical nature of the derived area law in terms of quantum information and Feynman contours.

Let the horizon surface consist of $N$ independent quantum binary degrees of freedom (qubits) generated by topological punctures of the closed contours of the phase fiber bundle $\mathcal{P}(\mathcal{M}, \mathcal{G})$. Each elementary puncture possesses two degenerate phase orientations ($d = 2$, binary qubit: $|\uparrow\rangle, |\downarrow\rangle$).

The total statistical number of microstates of the horizon surface $\Omega$ is:

\begin{equation}\label{eq:2.64}\label{eq:2.115}\tag{2.64}
\Omega = 2^N.
\end{equation}

According to Ludwig Boltzmann's fundamental formula, the entropy of the microscopic ensemble is:

\begin{equation}\label{eq:2.65}\label{eq:2.116}\tag{2.65}
S = k_B \ln\Omega = k_B \ln(2^N) = k_B N \ln 2.
\end{equation}

Equating the microscopic Boltzmann entropy~\eqref{eq:2.65} to the macroscopic thermodynamic Bekenstein~\cite{ref19}--Hawking~\cite{ref18} entropy~\eqref{eq:2.63}:

\begin{equation}\label{eq:2.66}\label{eq:2.117}\tag{2.66}
k_B N \ln 2 \equiv \frac{k_B A}{4 l_P^2}.
\end{equation}

From the fundamental identity~\eqref{eq:2.66}, we determine the elementary horizon surface area $a_{\text{bit}}$ corresponding strictly to one bit of quantum information:

\begin{equation}\label{eq:2.67}\label{eq:2.118}\tag{2.67}
a_{\text{bit}} \equiv \frac{A}{N} = 4\ln 2 \cdot l_P^2 = 4\ln 2 \, \frac{\hbar G}{c^3} \approx 2.7725887 \, l_P^2.
\end{equation}

The quantity $a_{\text{bit}} = 4\ln 2 \cdot l_P^2$ is a \textbf{theoretically proven consequence} of the mutual consistency between Hawking thermodynamics and quantum information science, rather than an ad hoc initial postulate. It defines the exact holographic information density: a black hole horizon encodes precisely 1 bit of information per every $4\ln 2$ Planck areas.

\subsubsection{Logarithmic Quantum Corrections}\label{sec:2.6.3}

In higher orders of quantum phase fluctuations around the saddle point, the Van Vleck-Morette determinant of second variations on the event horizon generates universal logarithmic corrections to the entropy:

\begin{equation}\label{eq:2.68}\label{eq:2.119}\tag{2.68}
S_{\text{quant}} = \frac{k_B A}{4 l_P^2} - \frac{3}{2} k_B \ln\left( \frac{A}{l_P^2} \right) + \mathcal{O}(1).
\end{equation}

The coefficient $-3/2$ is a topological invariant of conformal quantum zero modes on the horizon 2-sphere and does not violate the unitarity of quantum evolution, confirming the strict self-consistency of the contour framework.

\subsubsection{Hawking Radiation and the Four Laws of Black Hole Thermodynamics}\label{sec:2.6.4}

The fundamental thermodynamic characteristics of the black hole horizon derived in Sec.~\ref{sec:2.6.1}---the surface gravity $\kappa_{\text{grav}}$~\eqref{eq:2.58} and the Hawking quantum thermal radiation temperature $T_H$~\eqref{eq:2.59}:

\begin{equation*}\label{eq:2.120}
T_H = \frac{\hbar \kappa_{\text{grav}}}{2\pi c k_B} = \frac{\hbar c^3}{8\pi G M k_B},
\end{equation*}

determine the thermodynamic equilibrium of the black hole as a black body with quantum emission temperature $T_H$.

Black hole thermodynamics is governed by the four fundamental laws of Bardeen-Carter-Hawking \cite{ref20}, which constitute a direct physical manifestation of the laws of classical phenomenological thermodynamics:

\begin{enumerate}
\item \textbf{Zeroth Law:} The surface gravity $\kappa_{\text{grav}}$ is strictly constant across the entire closed horizon surface of a stationary black hole (the analogue of thermal equilibrium and temperature uniformity, $T = \text{const}$).

\item \textbf{First Law (Conservation of Energy):}  
The differential of the total mass-energy of the black hole $E = M c^2$ satisfies the thermodynamic balance:

\begin{equation}\label{eq:2.70}\label{eq:2.121}\tag{2.70}
d(M c^2) = \frac{\kappa_{\text{grav}} c^2}{8\pi G} dA + \Omega_H dJ + \Phi_H dQ = T_H dS + \Omega_H dJ + \Phi_H dQ,
\end{equation}

where $M$ is the total mass of the black hole;  
$A$ is the 2-surface area of the event horizon;  
$S$ is the Bekenstein-Hawking entropy derived in Eq.~\eqref{eq:2.63} ($S = k_B A / 4l_P^2$);  
$\Omega_H$ is the angular velocity of the event horizon;  
$J$ is the angular momentum of the black hole;  
$\Phi_H$ is the electrostatic potential on the horizon;  
$Q$ is the electric charge of the black hole.

\item \textbf{Second Law (Generalized Bekenstein Entropy Increase Law):}  
The total entropy of a closed system of the Universe, equal to the sum of the ordinary matter entropy and the horizon entropy of all black holes, is a non-decreasing function in any physical processes:

\begin{equation}\label{eq:2.71}\label{eq:2.122}\tag{2.71}
\Delta S_{\text{total}} = \Delta S_{\text{matter}} + \Delta S_{\text{BH}} \ge 0.
\end{equation}

\item \textbf{Third Law:} The horizon surface gravity $\kappa_{\text{grav}}$ cannot be reduced to zero through any finite sequence of physical operations or interactions (the unattainability of absolute zero Hawking temperature, $T_H \to 0$).
\end{enumerate}

Thus, the entire thermodynamics of black holes is a direct macroscopic manifestation of the quantum phase equilibrium of Feynman contours on the event horizon.

\subsection{Conclusions of Chapter 2}\label{sec:2.7}

\begin{enumerate}
\item On the basis of dynamic non-linear self-interaction of a massless spin-2 gauge field in Minkowski vacuum (the Feynman-Deser theorem), an exact derivation of the Einstein field equations of General Relativity is achieved without postulating an a priori curved spacetime geometry.
\item It is established that in strong and moderate fields ($a \gg a_0$), the theory of phase dynamics exactly reduces to General Relativity, guaranteeing the automatic fulfillment of classical relativistic tests (perihelion precession of Mercury, solar light deflection, Shapiro time delay, and LIGO gravitational waves).
\item From non-Abelian Wilson loops, the classical Yang-Mills action is deduced, asymptotic freedom and infrared topological confinement are analytically proven, and the mass gap problem is resolved (Yang-Mills mass gap $\Delta \approx 1.6\text{ GeV}$). It is proven that 99\% of the mass of baryonic matter ($m_p \approx 938\text{ MeV}/c^2$) is generated by closed quark-gluon contours. On the basis of the Bern-Carrasco-Johansson (BCJ) duality, the connection between gravity and gauge fields is proven as the double copy $\text{Gravity} = (\text{YM})^2$.
\item Fundamental Theorem~2 on the wave-optical quantum nature of the Equivalence Principle is proven: the identity $m_{\text{in}} \equiv m_{\text{grav}}$ is a rigorous consequence of de Broglie matter wave refraction in the vacuum phase index gradient $n(r) = 1 + 2GM/(rc^2)$ generated by the slowing of the Margolus-Levitin quantum clock.
\item A microscopic information-theoretic derivation of black hole thermodynamics is established: partitioning the horizon into quantum bit-cells of area $a_{\text{bit}} = 4\ln 2 \cdot l_P^2$ uniquely generates the Bekenstein-Hawking entropy $S = A / (4 l_P^2)$ with the universal $- \frac{3}{2} k_B \ln(A/l_P^2)$ logarithmic quantum correction without free adjustable parameters.
\end{enumerate}

% ------------------------------------------------------------------------------
% Sequence Verification Table: Equations (2.76) through (2.134)
% In linear chapter indexing, Part B corresponds to equations (2.76)-(2.134):
% (2.76), (2.77), (2.78), (2.79), (2.80), (2.81), (2.82), (2.83), (2.84), (2.85),
% (2.86), (2.87), (2.88), (2.89), (2.90), (2.91), (2.92), (2.93), (2.94), (2.95),
% (2.96), (2.97), (2.98), (2.99), (2.100), (2.101), (2.102), (2.103), (2.104), (2.105),
% (2.106), (2.107), (2.108), (2.109), (2.110), (2.111), (2.112), (2.113), (2.114), (2.115),
% (2.116), (2.117), (2.118), (2.119), (2.120), (2.121), (2.122), (2.123), (2.124), (2.125),
% (2.126), (2.127), (2.128), (2.129), (2.130), (2.131), (2.132), (2.133), (2.134).
% ------------------------------------------------------------------------------
